下面记录验证证明中的等号条件，以用于下一节的唯一性分析。

\begin{corollary}[验证等号条件]
\label{cor:verification-equality}
在定理\ref{thm:verification}的条件下，固定 $x_0\in\Dcal$。
对任意满足 $J_{x_0}(q)<\infty$ 的可行控制 $q\in\Uad(x_0)$，记
\[
\tau_q:=\tau_{\Acal}(x_0,q).
\]
在 $(0,\tau_q)$ 上几乎处处定义
\begin{equation}
R_q(t):=\frac{\mathrm d}{\mathrm d t}U(x_q(t))+pcq(t).
\label{eq:trajectory-verification-residual}
\end{equation}
则 $\tau_q<\infty$，$R_q\in L^1(0,\tau_q)$，且 $R_q\ge0$ 几乎处处。
此外，
\begin{equation}
J_{x_0}(q)-V(x_0)
=
\int_0^{\tau_q}R_q(t)\dd t
+
\int_{\tau_q}^{\infty}pcq(t)\dd t.
\label{eq:verification-gap}
\end{equation}
因此，在上述有限成本可行控制类中，$q$ 最优当且仅当
\begin{equation}
\begin{cases}
R_q(t)=0 & \text{几乎处处于 }(0,\tau_q),\\
q(t)=0 & \text{几乎处处于 }(\tau_q,\infty).
\end{cases}
\label{eq:verification-equality-conditions}
\end{equation}
当 $x_0\in\Acal$ 时，取 $\tau_q=0$，并将第一个积分理解为零。
\end{corollary}

\begin{proof}
由引理\ref{lem:finite-entry}，$\tau_q<\infty$，且
$x_q(\tau_q)\in\Acal$，故 $U(x_q(\tau_q))=0$。
若 $\tau_q>0$，引理\ref{lem:trajectory-verification}保证
$U\circ x_q$ 在 $[0,\tau_q]$ 上绝对连续，且 $R_q\ge0$ 几乎处处。
绝对连续函数的导数可积，而 $q$ 有界，所以
$R_q\in L^1(0,\tau_q)$。
由微积分基本定理，
\begin{align*}
\int_0^{\tau_q}R_q(t)\dd t
&=U(x_q(\tau_q))-U(x_0)
  +\int_0^{\tau_q}pcq(t)\dd t\\
&=-U(x_0)+\int_0^{\tau_q}pcq(t)\dd t.
\end{align*}
将成本积分在 $\tau_q$ 处分开，并使用
定理\ref{thm:verification}中的 $U(x_0)=V(x_0)$，
即得 \eqref{eq:verification-gap}。
当 $\tau_q=0$ 时，$U(x_0)=V(x_0)=0$，该恒等式同样成立。

式\eqref{eq:verification-gap}右端两项均非负。
非负可积函数的积分为零，当且仅当该函数几乎处处为零；
又因 $pc>0$，第二项为零当且仅当 $q=0$ 几乎处处于
$(\tau_q,\infty)$。
结合最优性的定义，得到 \eqref{eq:verification-equality-conditions}。
\end{proof}

\begin{remark}[等号条件与唯一性]
\label{rem:verification-equality-uniqueness}
在 $U$ 可微的不安全内部，普通链式法则给出
\[
R_q(t)=H_U(x_q(t),q(t))
\]
几乎处处成立。
在状态约束边界上，仍使用引理\ref{lem:trajectory-verification}
所建立的沿轨道导数关系。
推论\ref{cor:verification-equality}刻画最优控制的等号条件，
但不单独蕴含最优控制唯一。
唯一性还需结合状态约束、切换曲线的横穿性质及状态方程解的唯一性。
\end{remark}
